\section{Related Work}
\label{sec:related}

\textbf{Dataset characterization.} Trajectory quality and state--action coverage relate offline datasets to learner performance~\cite{schweighofer2022dataset}; ERI~\cite{swazinna2021quality} and the Bellman--Wasserstein distance~\cite{asadulaev2026bwd} score dataset quality, and image representations of a dataset predict the return of a learner trained on it~\cite{mateos2024images}. These methods output a score or a return at one configuration, some with reference data from the true environment. We use ERI and BWD as adapted input features under our head and the author-code TQ/SACo as a privileged comparison.

\textbf{Offline sample complexity and the effective horizon.} Coverage of the relevant policy and action gaps governs offline sample complexity~\cite{rashidinejad2021bridging,li2024sample,wang2022gapoffline}. The effective horizon links the lookahead depth at which greedy actions become correct to the effort of interactive RL~\cite{laidlaw2023effective,laidlaw2024stochastic}. Our planning-risk descriptor keeps its depth-and-gap structure but replaces the true MDP by a posterior and random rollouts by the logging behavior, which voids the guarantees; we treat it as a feature to be validated.

\textbf{Learning curves and off-policy evaluation.} Learning-curve extrapolation predicts final performance from partial training runs, parametrically~\cite{domhan2015speeding} or with prior-data fitted networks~\cite{adriaensen2023efficient}; learning curves can be non-monotone~\cite{viering2023shape}, which motivates our sustained-success rule. Fitted Q evaluation estimates a policy's value from logged data~\cite{le2019batch,voloshin2021empirical}, and offline policy selection ranks trained policies~\cite{yang2022selection,tang2021selection}. Our strongest baseline combines both: it trains the learner on pilot prefixes, scores it by FQE on held-out folds and extrapolates.

\textbf{Censored forecasting and tabular models.} Right-censored outcomes are commonly modeled with accelerated failure time models~\cite{barnwal2022aft} or discrete-time classifiers~\cite{kim2026survival}; proper scores~\cite{gneiting2007strictly} evaluate the forecasts. Tabular foundation models such as TabPFN~\cite{hollmann2025accurate} predict in context without fitting; we use TabPFN v2 as a strong generic learner on the pilot statistics. Reusing validation data biases evaluation~\cite{cawley2010selection,dwork2015adaptive}, which is why development and final panels are kept apart.
